\chapter{关键 SQL 与测试摘录}

\section{选课事务函数摘录}

本系统通过 \code{select_course_tx()} 封装数据库层选课入口。函数在调用方事务中执行，先锁定同一学生同一教学班的既有选课记录，再锁定目标教学班，最终由 \code{trg_enrollment_guard} 完成完整性兜底。

\begin{lstlisting}[style=sqlstyle]
CREATE OR REPLACE FUNCTION select_course_tx(
    p_student_id VARCHAR,
    p_offering_id BIGINT
) RETURNS BIGINT AS $$
...
    PERFORM 1
    FROM course_offering
    WHERE offering_id = p_offering_id
    FOR UPDATE;

    INSERT INTO enrollment (student_id, offering_id, status, select_time)
    VALUES (p_student_id, p_offering_id, 'selected', CURRENT_TIMESTAMP)
    RETURNING enrollment_id INTO v_enrollment_id;
...
$$ LANGUAGE plpgsql;
\end{lstlisting}

\section{角色一致性触发器摘录}

角色档案一致性由通用函数和三个触发器共同实现。下面片段说明普通外键之外，系统还会检查账户角色是否与档案表类型匹配。

\begin{lstlisting}[style=sqlstyle]
CREATE OR REPLACE FUNCTION check_user_role_fn(
    p_user_id BIGINT,
    p_expected_role VARCHAR
) RETURNS void AS $$
...
    IF v_role <> p_expected_role THEN
        RAISE EXCEPTION 'User % role is %, expected %',
            p_user_id, v_role, p_expected_role;
    END IF;
END;
$$ LANGUAGE plpgsql;

CREATE TRIGGER trg_student_role_check
BEFORE INSERT OR UPDATE OF user_id ON student
FOR EACH ROW
EXECUTE PROCEDURE trg_profile_role_check_fn('student');
\end{lstlisting}

\section{选课完整性触发器摘录}

\code{trg_enrollment_guard} 不维护 \code{selected_count}，而是在锁定教学班后统计 \code{enrollment} 中的 \code{selected} 记录数。

\begin{lstlisting}[style=sqlstyle]
SELECT course_id, semester_id, max_capacity, status
  INTO v_course_id, v_semester_id, v_max_capacity, v_offering_status
FROM course_offering
WHERE offering_id = NEW.offering_id
FOR UPDATE;

SELECT COUNT(*) INTO v_selected_count
FROM enrollment e
WHERE e.offering_id = NEW.offering_id
  AND e.status = 'selected'
  AND (v_old_enrollment_id IS NULL OR e.enrollment_id <> v_old_enrollment_id);

IF v_selected_count >= v_max_capacity THEN
    RAISE EXCEPTION 'Course offering % is full', NEW.offering_id;
END IF;
\end{lstlisting}

同课跨班限制也由同一触发器完成。该规则需要同时查询 \code{enrollment} 与 \code{course_offering}，因此不适合只用单表唯一约束表达。

\begin{lstlisting}[style=sqlstyle]
SELECT COUNT(*) INTO v_same_course_count
FROM enrollment e
JOIN course_offering co_same
  ON co_same.offering_id = e.offering_id
WHERE e.student_id = NEW.student_id
  AND e.status = 'selected'
  AND co_same.semester_id = v_semester_id
  AND co_same.course_id = v_course_id
  AND co_same.offering_id <> NEW.offering_id;

IF v_same_course_count > 0 THEN
    RAISE EXCEPTION
        'Student % has already selected another offering for course %',
        NEW.student_id, v_course_id;
END IF;
\end{lstlisting}

\section{成绩审计触发器摘录}

成绩审计触发器要求业务层在同一事务中传入当前操作者，避免把日志操作者写死为管理员或空值。

\begin{lstlisting}[style=sqlstyle]
BEGIN
    v_changed_by_user_id := current_setting('app.current_user_id')::BIGINT;
EXCEPTION WHEN OTHERS THEN
    RAISE EXCEPTION
      'Score update requires app.current_user_id in the current transaction';
END;

INSERT INTO score_change_log
    (enrollment_id, old_score, new_score, changed_by_user_id, reason)
VALUES
    (NEW.enrollment_id, OLD.final_score, NEW.final_score,
     v_changed_by_user_id, v_reason);
\end{lstlisting}

\section{视图 SQL 摘录}

\code{v_offering_selected_count} 通过聚合实时得到已选人数和剩余容量，使基础表无需保存派生字段 \code{selected_count}。

\begin{lstlisting}[style=sqlstyle]
CREATE OR REPLACE VIEW v_offering_selected_count AS
SELECT
    co.offering_id,
    co.max_capacity,
    COALESCE(COUNT(e.enrollment_id), 0)::INTEGER AS selected_count,
    (co.max_capacity - COALESCE(COUNT(e.enrollment_id), 0))::INTEGER
        AS remaining_capacity
FROM course_offering co
LEFT JOIN enrollment e
  ON e.offering_id = co.offering_id
 AND e.status = 'selected'
GROUP BY co.offering_id, co.max_capacity;
\end{lstlisting}

\section{SQL 测试结果摘录}

\code{opengauss_setup/sql/test_triggers_constraints.sql} 使用事务和匿名块捕获预期异常，最后执行 \code{ROLLBACK}，不会污染样例数据。测试覆盖角色错配、先修自指、先修环路、容量超选、时间冲突、同课跨班、重复选课、成绩审计、退课限制和视图查询。一次测试输出摘录如下：

\begin{lstlisting}[style=codestyle]
PASS role check teacher-as-student failed
PASS self prerequisite failed
PASS prerequisite cycle failed
PASS capacity full failed
PASS timetable conflict failed
PASS same course different offering failed
PASS unique student/offering failed
PASS score update without user failed
PASS score audit trigger inserted log row
PASS completed enrollment drop failed
Trigger and constraint tests finished with rollback.
\end{lstlisting}

\section{并发抢课测试脚本摘录}

\code{scripts/concurrency_enrollment_test.py} 用于在两个独立数据库连接中同时调用 \code{select_course_tx()}。并发测试在本地 openGauss 容器中构造容量为 1 的教学班，最终只有一个学生选课成功。测试输出摘录如下：

\begin{lstlisting}[style=codestyle]
TST_CONC_STU_1: FAILED - Course offering 42 is full
TST_CONC_STU_2: SUCCESS - selected
success_count = 1
failed_count = 1
selected_count_in_db = 1
\end{lstlisting}

\section{数据库角色授权脚本摘录}

\code{grant_roles_20260608.sql} 体现数据库管理中的最小权限思想。本地生成的 \code{opengauss_setup/sql/local/grant_demo_db_accounts_20260611.sql} 用于演示数据库登录账号与角色绑定。由于本地绑定脚本包含数据库账号密码，不进入公开仓库和报告正文。

\begin{lstlisting}[style=sqlstyle]
CREATE ROLE db_student_role;
CREATE ROLE db_teacher_role;
CREATE ROLE db_admin_role;
CREATE ROLE db_app_role;

REVOKE ALL ON ALL TABLES IN SCHEMA public FROM PUBLIC;

GRANT SELECT ON
    v_course_offering_detail,
    v_offering_selected_count
TO db_student_role;

GRANT EXECUTE ON FUNCTION select_course_tx(VARCHAR, BIGINT)
TO db_app_role;
\end{lstlisting}

\section{演示账号凭据策略摘录}

演示业务账号使用业务身份编号作为登录名：学生用学号，教师用教师号，管理员用管理员编号。明文初始密码只保存在本地 \code{secrets/DEMO_ACCOUNT_CREDENTIALS.md}，SQL 中只保存 SHA-256 哈希。

\begin{lstlisting}[style=codestyle]
student username = student.student_id      # e.g. 20240001
teacher username = teacher.teacher_id      # e.g. T001
admin username   = admin_profile.admin_id  # e.g. A001

python3 scripts/generate_large_demo_dataset.py
\end{lstlisting}

\section{绩点规则版本化 SQL 摘录}

绩点不作为 \code{enrollment} 的基础字段保存，而是由规则版本和成绩区间计算。下面摘录展示 \code{grade_policy}、\code{grade_scale} 和计算函数的核心结构。

\begin{lstlisting}[style=sqlstyle]
CREATE TABLE grade_policy (
    policy_id BIGSERIAL PRIMARY KEY,
    policy_code VARCHAR(30) NOT NULL,
    policy_name VARCHAR(100) NOT NULL,
    version_no VARCHAR(20) NOT NULL,
    effective_from DATE NOT NULL,
    effective_to DATE,
    status VARCHAR(20) NOT NULL DEFAULT 'active',
    UNIQUE (policy_code, version_no)
);

CREATE TABLE grade_scale (
    scale_id BIGSERIAL PRIMARY KEY,
    policy_id BIGINT NOT NULL REFERENCES grade_policy(policy_id),
    min_score DECIMAL(5,2) NOT NULL,
    max_score DECIMAL(5,2) NOT NULL,
    gpa_point DECIMAL(3,2) NOT NULL,
    grade_label VARCHAR(10)
);
\end{lstlisting}

\begin{lstlisting}[style=sqlstyle]
CREATE OR REPLACE FUNCTION calculate_gpa(
    p_score DECIMAL,
    p_policy_id BIGINT
) RETURNS DECIMAL(3,2) AS $$
...
    SELECT COUNT(*), MIN(gpa_point)
      INTO v_match_count, v_gpa
    FROM grade_scale
    WHERE policy_id = p_policy_id
      AND p_score >= min_score
      AND p_score <= max_score;
...
$$ LANGUAGE plpgsql;
\end{lstlisting}

\code{test_grade_policy.sql} 覆盖绩点规则插入、区间重叠失败、\code{calculate_gpa(95)}、\code{calculate_gpa(59)} 以及绩点视图查询。回滚型测试输出摘录如下：

\begin{lstlisting}[style=codestyle]
PASS calculate_gpa 95 -> 4.00
PASS calculate_gpa 59 -> 0.00
PASS overlap grade scale failed
PASS v_enrollment_grade_detail returned ... GPA rows
PASS v_student_gpa_summary returned ... summary rows
Grade policy tests finished with rollback.
\end{lstlisting}

\section{大规模演示数据验证摘录}

最终演示数据由 \code{scripts/generate_large_demo_dataset.py} 使用固定随机种子生成，导入脚本位于本地 \code{opengauss_setup/sql/local/seed_large_demo_dataset_20260611.sql}。明文初始密码只保存在本地 \code{secrets/DEMO_ACCOUNT_CREDENTIALS.md}，不进入报告正文。数据规模摘要如下：

\begin{lstlisting}[style=codestyle]
Departments: 4
Students: 100
Teachers: 30
Courses: 50
Course offerings: 175
Schedule slots: 214
Enrollments: 1480
Completed enrollments: 1000
Selected enrollments: 425
Dropped enrollments: 55
Score change logs: 80
\end{lstlisting}

\code{validate_large_demo_dataset.sql} 对大规模数据进行只读验证，检查教学班容量、教室容量、重复选课、同课跨班、时间冲突、先修课、成绩状态、账号角色、先修环路和成绩审计日志。关键检查片段如下：

\begin{lstlisting}[style=sqlstyle]
WITH selected_count AS (
    SELECT offering_id, COUNT(*) AS selected_count
    FROM enrollment
    WHERE status = 'selected'
    GROUP BY offering_id
)
SELECT COUNT(*) AS violation_count
FROM course_offering co
LEFT JOIN selected_count sc ON sc.offering_id = co.offering_id
WHERE COALESCE(sc.selected_count, 0) > co.max_capacity;

SELECT COUNT(*) AS violation_count
FROM enrollment e1
JOIN enrollment e2
  ON e1.student_id = e2.student_id
 AND e1.enrollment_id < e2.enrollment_id
JOIN course_offering co1 ON co1.offering_id = e1.offering_id
JOIN course_offering co2 ON co2.offering_id = e2.offering_id
JOIN course_schedule cs1 ON cs1.offering_id = co1.offering_id
JOIN course_schedule cs2 ON cs2.offering_id = co2.offering_id
WHERE e1.status = 'selected'
  AND e2.status = 'selected'
  AND co1.semester_id = co2.semester_id
  AND cs1.weekday = cs2.weekday
  AND cs1.start_section <= cs2.end_section
  AND cs1.end_section >= cs2.start_section;
\end{lstlisting}

本地 openGauss 验证结果显示，容量超限、教室容量超限、重复选课、同课跨班、时间冲突、先修课不满足、成绩状态异常、角色错配、先修自指或环路、审计日志无效引用等违规项均为 0。绩点视图 smoke check 返回 \code{v_enrollment_grade_detail} 1480 行、\code{v_student_gpa_summary} 284 行。

\section{教师评价扩展设计}

教师评价不作为本系统当前主线功能实现。若后续扩展，可新增 \code{teaching_evaluation} 表，并把它作为选课记录完成后的附属业务实体。

\begin{lstlisting}[style=sqlstyle]
teaching_evaluation(
    evaluation_id,
    enrollment_id,
    rating,
    comment,
    submitted_at
)
\end{lstlisting}

设计上应保证一个 \code{enrollment} 最多评价一次，评分范围为 1 到 5，且只有 \code{completed} 状态的选课记录可以评价。教师侧宜查看聚合结果，例如平均分与评价人数，而不是直接查看学生身份。
